% Part: history
% Chapter: biographies 
% Section: georg-cantor

\documentclass[../../../include/open-logic-section]{subfiles}

\begin{document}

\olfileid{his}{bio}{can}

\olsection{گیورگ کانتور}

\olphoto{cantor-georg}{گیورگ کانتور}

گیورگ کانتور (\textsc{گے} اورگ \textsc{کان} تور) دی اک مُڈھلی سوانح عمری وچ دعویٰ کیتا گیا سی کہ اوہ روس دے سینٹ پیٹرزبرگ ول جاندے اک بحری جہاز وچ جمیا تے اوتھے ای لبھیا سی، تے اوہدے ماں پیو دا پتا نہیں سی۔ پر ایہ گل سچ نہیں۔ اوہ 1845 وچ سینٹ پیٹرزبرگ وچ ضرور جمیا سی۔

کانتور نے 1867 وچ برلن یونیورسٹی توں ریاضی وچ ڈاکٹریٹ لئی۔ اوہ سیٹ نظریے وچ اپنے کم لئی جانیا جاندا اے، تے سیٹ نظریے نوں تحقیق دے اک وکھرے شعبے وجوں قائم کرن دا سہرا اوہدے سر اے۔ اوہ پہلا بندہ سی جیہنے ثابت کیتا کہ وکھ وکھ وڈائی دے لامتناہی سیٹ ہُندے نیں۔ اوہدے نظریاں، خاص کر کے لامتناہیت دے نظریے، نے اوس ویلے دے ریاضی داناں وچ بڑی بحث چھیڑی، تے اوہدا کم اختلاف دا باعث رہیا۔

کانتور دے مذہبی یقین تے اوہدا ریاضیاتی کم اینے جڑے سن کہ اوہناں نوں وکھ نہیں کیتا جا سکدا۔ اوہنے ایہ دعویٰ وی کیتا کہ لامتناہی درجیاں والے عدداں دا نظریہ رب نے اوہنوں سِدھا دسیا سی۔ پچھلی عمر وچ کانتور ذہنی بیماری دا شکار رہیا۔ 1894 توں اوہ ہسپتال داخل ہون لگ پیا، تے عمر دے آخری سالاں وچ ایہ ہور واری واری ہویا۔ اوہدے کم اُتے سخت تنقید، جیہدے وچ ریاضی دان لیوپولڈ کرونیکر نال ان بن وی شامل سی، اوہدی افسردگی تے ریاضی وچ دلچسپی گھٹ جان دا سبب بنی۔ افسردگی دے دوراں وچ کانتور فلسفے تے ادب ول مُڑ جاندا سی۔ اوہنے ایہ نظریہ وی چھپوایا کہ شیکسپیئر دے ڈرامے فرانسس بیکن نے لکھے سن۔

کانتور 6 جنوری 1918 نوں ہالے دے اک علاج گھر وچ وفات پا گیا۔

\begin{reading}
کانتور دیاں پوریاں سوانح عمریاں لئی \citet{Dauben1990} تے \citet{Grattan-Guinness1971} ویکھو۔ کانتور دے رواج توں ہٹ کے خیالات دا بیان بی بی سی ریڈیو 4 دے پروگرام \emph{A Brief History of
  Mathematics} وچ وی ملدا اے \citep{Sautoy2014}۔ جے تسی کانتور دے نظریے ریپ دی صورت وچ سُننا چاہندے او، تاں \citet{Rose2012} ویکھو۔
\end{reading}

\end{document}
